% Pluricanonical character/decomposition and Artin conductor degrees for the S5-regular closure cover.

\begin{theorem}[\texorpdfstring{$(5)$}{(5)}-循环固定点的切表示全谱性]\label{thm:xi-terminal-zm-delta-s5-5cycle-tangent-spectrum}
\leanverified{paper\_xi\_terminal\_zm\_delta\_s5\_5cycle\_tangent\_spectrum}
令 \(\widehat\pi_\delta:X\to\PP^1_\delta\) 为结论 \ref{con:xi-terminal-zm-delta-s5-regular-closure-genus-109} 中的正则 \(S_5\)-伽罗瓦覆盖，并记 \(G:=S_5\)。
取任一 \(5\)-循环 \(\sigma\in G\)，记其固定点集为 \(X^\sigma\)。
则 \(X^\sigma\) 非空，且事实上 \(|X^\sigma|=4\)。
并且对每个固定点 \(P\in X^\sigma\)，导数作用
\[
d\sigma_P:T_PX\to T_PX
\]
之本征值为某个非平凡五次单位根；按固定点计重后，切本征值的多重集合恰为
\[
\left\{\zeta_5,\zeta_5^{2},\zeta_5^{3},\zeta_5^{4}\right\},
\qquad \zeta_5:=\exp(2\pi i/5).
\]
\end{theorem}
\begin{proof}
由护照 \((5)\cdot(2)^6\)（定理 \ref{thm:xi-terminal-zm-delta-map-hurwitz-passport-s5}），\(\widehat\pi_\delta\) 的非平凡惯性子群仅可能为阶 \(2\) 或阶 \(5\) 的循环群。
因此若 \(g\in G\) 固定某点 \(P\in X\)，则 \(g\) 必属于某个分歧点的惯性群的共轭，特别地，若 \(g\) 为 \(5\)-循环，则必在某个阶 \(5\) 的惯性群上出现。
\par
取无穷远分歧值 \(\infty\in\PP^1_\delta\)（其惯性阶为 \(5\)；参见定理 \ref{thm:xi-terminal-zm-delta-map-hurwitz-passport-s5}），并取其上方一点 \(P_0\in X\) 使稳定子为阶 \(5\) 的惯性群
\(
G_{P_0}=\langle\sigma_0\rangle
\)
（其中 \(\sigma_0\) 为某个 \(5\)-循环）。
由于 \(S_5\) 中所有 \(5\)-循环共轭，存在 \(g_0\in G\) 使得 \(\sigma=g_0\sigma_0g_0^{-1}\)，令 \(P:=g_0P_0\)，则 \(G_P=\langle\sigma\rangle\) 且 \(P\in X^\sigma\)，从而 \(X^\sigma\neq\varnothing\)。
\par
下面给出固定点数。
记 \(H_0:=\langle\sigma_0\rangle\)。
由于 \(\widehat\pi_\delta\) 为正则 \(G\)-覆盖，\(\widehat\pi_\delta^{-1}(\infty)\) 可识别为陪集空间 \(G/H_0\)。
\(\sigma\) 在 \(G/H_0\) 上的固定点个数为
\[
\#\{gH_0\in G/H_0:\ \sigma\cdot gH_0=gH_0\}
=\frac{\#\{g\in G:\ g^{-1}\sigma g\in H_0\}}{|H_0|}.
\]
在 \(S_5\) 中，一个 \(5\)-循环生成其唯一的 Sylow-\(5\) 子群，故条件 \(g^{-1}\sigma g\in H_0\) 等价于 \(gH_0g^{-1}=\langle\sigma\rangle\)。
取任意 \(g_0\) 满足 \(g_0H_0g_0^{-1}=\langle\sigma\rangle\)，则所有满足该条件的 \(g\) 组成右陪集 \(g_0N_G(H_0)\)，其大小为 \(|N_G(H_0)|\)。
而对 \(H_0\simeq C_5\)，有
\(
N_G(H_0)/H_0\cong \Aut(H_0)\cong(\ZZ/5)^\times
\)
且 \(|(\ZZ/5)^\times|=4\)，因此 \(|N_G(H_0)|=20\)，从而固定点个数为 \(20/5=4\)。
这给出 \(|X^\sigma|=4\)。
\par
最后确定切本征值。
在 \(P_0\) 的完备局部环上，\(\sigma_0\) 的作用为忠实阶 \(5\) 的解析旋转，故存在局部参数 \(t\) 使 \(\sigma_0(t)=\zeta_5 t\)，于是 \(d\sigma_0|_{P_0}\) 的本征值为 \(\zeta_5\)。
对任意 \(h\in N_G(H_0)\)，存在唯一 \(a(h)\in(\ZZ/5)^\times\) 使 \(h^{-1}\sigma_0h=\sigma_0^{a(h)}\)。
记 \(P_h:=hg_0^{-1}P\in \widehat\pi_\delta^{-1}(\infty)\) 并取局部参数 \(t_h:=t\circ h^{-1}\)，则
\(
\sigma_0(t_h)=\zeta_5^{a(h)}t_h
\)。
当 \(h\) 在 \(N_G(H_0)/H_0\simeq(\ZZ/5)^\times\) 中遍历时，\(a(h)\) 遍历 \(\{1,2,3,4\}\)。
由上述对固定点计数的构造，\(\sigma\) 的四个固定点可由这些 \(h\) 的不同像代表得到，从而切本征值多重集合为 \(\{\zeta_5^k\}_{k=1}^4\)。
证毕。
\end{proof}

\begin{theorem}[多重典范表示的三支撑特征函数与周期性]\label{thm:xi-terminal-zm-delta-s5-pluricanonical-character-support}
\leanverified{paper\_xi\_terminal\_zm\_delta\_s5\_pluricanonical\_character\_support}
令 \(K_X\) 为 \(X\) 的典范线丛，并对整数 \(n\ge 2\) 定义
\[
V_n:=H^0(X,nK_X).
\]
令 \(\chi_n\) 为 \(G=S_5\) 在 \(V_n\) 上的特征函数。
则成立：
\begin{enumerate}
  \item 对任意既非换位亦非 \(5\)-循环且非恒等元的 \(g\in S_5\)，有 \(\chi_n(g)=0\)。
  \item 若 \(\tau\) 为任一换位，则
  \[
  \chi_n(\tau)=18(-1)^n.
  \]
  \item 若 \(\sigma\) 为任一 \(5\)-循环，则
  \[
  \chi_n(\sigma)=u_n,
  \qquad
  u_n:=\sum_{k=1}^4\frac{\zeta_5^{kn}}{1-\zeta_5^{k}}.
  \]
  因而 \(u_n\) 仅依赖于 \(n\bmod 5\)，并满足显式取值
  \[
  u_n=
  \begin{cases}
  2,&n\equiv 0\pmod 5,\\
  -2,&n\equiv 1\pmod 5,\\
  -1,&n\equiv 2\pmod 5,\\
  0,&n\equiv 3\pmod 5,\\
  1,&n\equiv 4\pmod 5.
  \end{cases}
  \]
  \item 恒等元处
  \[
  \chi_n(1)=\dim_\CC V_n=108(2n-1).
  \]
\end{enumerate}
\end{theorem}
\begin{proof}
\textbf{(1)} 若 \(g\in S_5\) 固定某点 \(P\in X\)，则 \(g\in G_P\)。
而由护照 \((5)\cdot(2)^6\) 可知 \(G_P\) 仅可能为阶 \(2\) 或阶 \(5\) 的循环群，故非平凡固定元必为换位或 \(5\)-循环。
\par
\textbf{(4)} 由 \(g(X)=109\)，对 \(n\ge 2\) 有 \(\deg(nK_X)=n(2g-2)=216n\)，并且 \(H^1(X,nK_X)=0\)。
由 Riemann--Roch 得
\[
\dim V_n=h^0(X,nK_X)=\deg(nK_X)+1-g=216n+1-109=108(2n-1).
\]
\par
\textbf{(2)(3)} 对任意自同构 \(g\) 与 \(n\ge 2\)，由于 \(H^1(X,nK_X)=0\)，全纯 Lefschetz 固定点公式给出
\[
\chi_n(g)=\mathrm{tr}\!\left(g\mid H^0(X,nK_X)\right)
=\sum_{P\in X^g}\frac{(dg_P)^n}{1-dg_P},
\]
其中 \(dg_P\) 为 \(g\) 在切空间上的导数本征值。
\par
若 \(g=\tau\) 为换位，则每个固定点处 \(dg_P=-1\)，从而
\[
\chi_n(\tau)=\sum_{P\in X^\tau}\frac{(-1)^n}{1-(-1)}
=\frac{(-1)^n}{2}\cdot |X^\tau|.
\]
为确定 \(|X^\tau|\)，取 \(n=1\) 并用 \(H^0(X,K_X)\cong H^0(X,\Omega_X^1)\) 的分解（结论 \ref{con:xi-terminal-zm-delta-s5-holomorphic-oneforms-chevalley-weil}）。
对换位 \(\tau\) 的迹为
\[
\mathrm{tr}\!\left(\tau\mid H^0(X,K_X)\right)
=2\cdot(-1)+1\cdot 2+3\cdot 1+5\cdot 0+6\cdot(-1)+7\cdot(-2)
=-17,
\]
其中用到 \(S_5\) 不可约特征在换位处的标准取值
\(
\chi_\varepsilon(\tau)=-1,\ \chi_{(4,1)}(\tau)=2,\ \chi_{(3,2)}(\tau)=1,\ \chi_{(3,1,1)}(\tau)=0,\ \chi_{(2,2,1)}(\tau)=-1,\ \chi_{(2,1,1,1)}(\tau)=-2
\)。
另一方面，\(H^1(X,K_X)\cong H^0(X,\mathcal O_X)^\vee\cong\CC\) 且 \(\tau\) 作用为恒等，故 Lefschetz 公式给出
\[
-17-1=\sum_{P\in X^\tau}\frac{-1}{1-(-1)}=-\frac{|X^\tau|}{2},
\]
从而 \(|X^\tau|/2=18\)，代回即得 \(\chi_n(\tau)=18(-1)^n\)。
\par
若 \(g=\sigma\) 为 \(5\)-循环，则由定理 \ref{thm:xi-terminal-zm-delta-s5-5cycle-tangent-spectrum}，固定点为四个且切本征值为 \(\{\zeta_5^k\}_{k=1}^4\)，故
\[
\chi_n(\sigma)=\sum_{k=1}^4\frac{\zeta_5^{kn}}{1-\zeta_5^k}=u_n.
\]
周期性与取值表由直接化简得到。
证毕。
\end{proof}

\begin{theorem}[\texorpdfstring{$H^0(X,nK_X)$}{H0(X,nKX)} 的 \texorpdfstring{$S_5$}{S5} 不可约分解：显式重数与周期 \(10\)]\label{thm:xi-terminal-zm-delta-s5-pluricanonical-irrep-decomposition-period10}
设 \(n\ge 2\)。
用分拆记号表示 \(S_5\) 的七个不可约复表示：
\[
(5),\ (4,1),\ (3,2),\ (3,1,1),\ (2,2,1),\ (2,1,1,1),\ (1^5),
\]
并记对应表示为 \(V_\lambda\)，其中 \(V_{(1^5)}=\varepsilon\) 为符号表示。
置
\[
s_n:=(-1)^n,\qquad
u_n:=\sum_{k=1}^4\frac{\zeta_5^{kn}}{1-\zeta_5^{k}}.
\]
则在 \(S_5\)-表示意义下存在唯一分解
\[
V_n\ \cong\ \bigoplus_{\lambda\vdash 5} V_\lambda^{\oplus m_\lambda(n)},
\]
其中各重数给出为
\[
\begin{aligned}
m_{(5)}(n)&=\frac{18n-9+15s_n+2u_n}{10},&
m_{(1^5)}(n)&=\frac{18n-9-15s_n+2u_n}{10},\\[2mm]
m_{(4,1)}(n)&=\frac{36n-18+15s_n-u_n}{5},&
m_{(2,1,1,1)}(n)&=\frac{36n-18-15s_n-u_n}{5},\\[2mm]
m_{(3,2)}(n)&=\frac{18n-9+3s_n}{2},&
m_{(2,2,1)}(n)&=\frac{18n-9-3s_n}{2},\\[2mm]
m_{(3,1,1)}(n)&=\frac{54n-27+u_n}{5}.&
&
\end{aligned}
\]
并且 \(m_\lambda(n)\in\ZZ_{\ge 0}\)；其 \(n\)-依赖为模 \(10\) 的拟多项式。
\end{theorem}
\begin{proof}
令 \(\chi_\lambda\) 为 \(V_\lambda\) 的不可约特征。
由特征正交性，
\[
m_\lambda(n)=\langle \chi_n,\chi_\lambda\rangle
=\frac{1}{|S_5|}\sum_{g\in S_5}\chi_n(g)\chi_\lambda(g).
\]
由定理 \ref{thm:xi-terminal-zm-delta-s5-pluricanonical-character-support}，\(\chi_n\) 仅在三类共轭类上取非零值：恒等元、换位类与 \(5\)-循环类。
相应共轭类大小分别为 \(1,10,24\)，故
\[
m_\lambda(n)
=\frac{1}{120}\Bigl(\chi_n(1)\chi_\lambda(1)+10\,\chi_n(\tau)\chi_\lambda(\tau)+24\,\chi_n(\sigma)\chi_\lambda(\sigma)\Bigr),
\]
其中 \(\tau\) 为换位、\(\sigma\) 为 \(5\)-循环。
代入
\(
\chi_n(1)=108(2n-1),\ \chi_n(\tau)=18s_n,\ \chi_n(\sigma)=u_n
\)
并使用 \(S_5\) 不可约特征在 \(\tau,\sigma\) 处的标准取值（可由 Murnaghan--Nakayama 规则得到），逐一化简即可得到所列闭式。
整数性与非负性由表示论结构保证。
证毕。
\end{proof}

\begin{proposition}\label{prop:xi-terminal-zm-delta-s5-regular-closure-gonality-lowerbound-5}
沿用结论 \ref{con:xi-terminal-zm-delta-s5-regular-closure-genus-109} 的记号，则
\[
\operatorname{gon}(X)\ge 5.
\]
\end{proposition}
\begin{proof}
取点稳定子 \(H\le S_5\)（故 \(H\cong S_4\) 且 \([S_5:H]=5\)）。
由结论 \ref{con:xi-terminal-zm-delta-s5-regular-closure-genus-109}，\(L^{H}=K=\QQ(E)\)，从而商曲线 \(X/H\) 与 \(E\) 双有理同构，特别地 \(g(X/H)=1\)。
令
\(
\pi_H:X\to X/H
\)
为自然商映射，则 \(\deg(\pi_H)=|H|=24\)。
\par
反设存在非常值映射 \(f:X\to\PP^1\) 且 \(\deg(f)=d\le 4\)。
若 \(f\) 经由 \(\pi_H\) 分解，则 \(\deg(f)=\deg(\pi_H)\deg(\bar f)\) 必为 \(24\) 的倍数，矛盾；
故 \(f\) 不经由 \(\pi_H\) 分解。
由 Castelnuovo--Severi 不等式（应用于两条不同的 pencil \(\pi_H\) 与 \(f\)）得
\[
g(X)\le 24\,g(X/H)+d\,g(\PP^1)+(24-1)(d-1)=24+23(d-1)=23d+1.
\]
当 \(d\le 4\) 时右端至多为 \(93\)，与 \(g(X)=109\) 矛盾。
故 \(\operatorname{gon}(X)\ge 5\)。证毕。
\end{proof}

\begin{theorem}\label{thm:xi-terminal-zm-delta-s5-canonical-quadrics-irrep-decomposition}
沿用结论 \ref{con:xi-terminal-zm-delta-s5-regular-closure-genus-109} 的记号。
令
\[
V:=H^0(X,K_X)\cong H^0(X,\Omega_X^1),
\qquad
\mu_2:\mathrm{Sym}^2V\longrightarrow H^0(X,2K_X)
\]
为乘法映射，并定义二次关系空间 \(Q_X:=\ker(\mu_2)\)。
则 \(Q_X\) 作为 \(S_5\)-模具有如下完全不可约分解：
\[
Q_X\ \cong\
V_{(5)}^{\oplus 70}\ \oplus\ V_{(1^5)}^{\oplus 39}\ \oplus\ V_{(4,1)}^{\oplus 219}\ \oplus\ V_{(2,1,1,1)}^{\oplus 159}\
\oplus\ V_{(3,2)}^{\oplus 258}\ \oplus\ V_{(2,2,1)}^{\oplus 228}\ \oplus\ V_{(3,1,1)}^{\oplus 270}.
\]
从而 \(\dim_\CC Q_X=5671\)，并且
\[
\dim_\CC(Q_X^{S_5})=70.
\]
此外，由命题 \ref{prop:xi-terminal-zm-delta-s5-regular-closure-gonality-lowerbound-5} 得 \(\operatorname{gon}(X)\ge 5\)，因而 \(X\) 非超椭圆且非三合页；
由 Petri 定理可知 \(X\) 的典范理想由 \(Q_X\) 所对应的二次式生成。
\end{theorem}
\begin{proof}
由结论 \ref{con:xi-terminal-zm-delta-s5-holomorphic-oneforms-chevalley-weil} 得到 \(V\) 的 \(S_5\)-等型分解，从而可计算其特征函数 \(\chi_V\)。
由表示论恒等式
\[
\chi_{\mathrm{Sym}^2V}(g)=\frac{1}{2}\bigl(\chi_V(g)^2+\chi_V(g^2)\bigr)
\]
得到 \(\mathrm{Sym}^2V\) 的特征。
另一方面，定理 \ref{thm:xi-terminal-zm-delta-s5-pluricanonical-character-support} 在 \(n=2\) 时给出 \(H^0(X,2K_X)\) 的特征函数 \(\chi_2\)。
因此
\[
\chi_{Q_X}=\chi_{\mathrm{Sym}^2V}-\chi_2.
\]
将 \(\chi_{Q_X}\) 与 \(S_5\) 的不可约特征作内积 \(m_\lambda=\langle \chi_{Q_X},\chi_\lambda\rangle\)，即可得到上式所列各不可约分量的重数，从而得到完全分解。
\par
维数断言由 Riemann--Roch 给出 \(\dim H^0(X,2K_X)=3g(X)-3=324\) 与 \(\dim \mathrm{Sym}^2V=g(X)(g(X)+1)/2=5995\)，故 \(\dim Q_X=5995-324=5671\)。
不变维数为平凡表示的重数 \(m_{(5)}(Q_X)=70\)。
最后一条由命题 \ref{prop:xi-terminal-zm-delta-s5-regular-closure-gonality-lowerbound-5} 排除三合页并结合 \(g(X)=109\) 排除平面五次的特殊情形，应用 Petri 定理得到典范理想由二次式生成。证毕。
\end{proof}

\begin{proposition}\label{prop:xi-terminal-zm-delta-s5-equivariant-deformation-dimension-4}
沿用结论 \ref{con:xi-terminal-zm-delta-s5-regular-closure-genus-109} 的记号，则保持 \(S_5\)-作用的一阶等变变形空间满足
\[
\dim_\CC H^1(T_X)^{S_5}=4.
\]
因此该护照 \((5)\cdot(2)^6\) 的 Hurwitz 参数空间在 \(X\) 处的切空间维数亦为 \(4\)（分支点数为 \(7\)，故自由度为 \(7-3\)）。
\end{proposition}
\begin{proof}
由 Serre 对偶 \(H^1(T_X)\cong H^0(X,2K_X)^\vee\) 且该同构与 \(S_5\)-作用相容，故
\[
\dim H^1(T_X)^{S_5}=\dim H^0(X,2K_X)^{S_5}.
\]
由定理 \ref{thm:xi-terminal-zm-delta-s5-pluricanonical-irrep-decomposition-period10} 取 \(n=2\)，其中 \(s_2=1\)、\(u_2=-1\)，得
\[
\dim H^0(X,2K_X)^{S_5}=m_{(5)}(2)=\frac{18\cdot 2-9+15\cdot 1+2\cdot(-1)}{10}=4.
\]
最后，对 \(\PP^1\) 上 \(7\) 个有序分支点的配置空间取 \(\mathrm{PGL}_2\) 商得到维数 \(7-3=4\)，即 Hurwitz 参数空间的切空间维数。证毕。
\end{proof}

\begin{proposition}[不可约 Artin 局部系统的全局导子度数向量]\label{prop:xi-terminal-zm-delta-s5-artin-conductor-degrees}
令 \(K=\QQ(\delta)\)、\(L=\QQ(X)\)，则 \(L/K\) 为 \(S_5\)-伽罗瓦扩张。
对每个不可约表示 \(V_\lambda\) 取对应的 Artin 表示 \(\rho_\lambda\)（例如取任意素数 \(\ell\ne 2,5\) 的 \(\ell\)-进实现），并记 \(\mathfrak f(\rho_\lambda)\) 为其 Artin 导子除子。
则
\[
\deg\mathfrak f(\rho_\lambda)
=6\bigl(\dim V_\lambda-\dim(V_\lambda^{C_2})\bigr)
+\bigl(\dim V_\lambda-\dim(V_\lambda^{C_5})\bigr),
\]
其中 \(C_2\)（分别 \(C_5\)）为由换位（分别 \(5\)-循环）生成的惯性循环群。
由结论 \ref{con:xi-terminal-zm-delta-s5-holomorphic-oneforms-chevalley-weil} 中给出的 \(\dim(V_\lambda^{C_2})\)、\(\dim(V_\lambda^{C_5})\) 数据可得（平凡表示 \((5)\) 的导子为 \(0\)）：
\[
\begin{array}{c|cccccc}
\lambda & (1^5) & (4,1) & (3,2) & (3,1,1) & (2,2,1) & (2,1,1,1)\\ \hline
\deg\mathfrak f(\rho_\lambda) & 6 & 10 & 16 & 22 & 22 & 22
\end{array}
\]
\end{proposition}
\begin{proof}
本覆盖的分歧惯性共轭类为一个 \(C_5\) 分歧点与六个 \(C_2\) 分歧点（护照 \((5)\cdot(2)^6\)）。
并且分歧为 tame，故在每个分歧点 \(v\) 处的 Artin 导子指数满足
\(
a_v(\rho)=\dim\rho-\dim(\rho^{I_v})
\)。
对全部分歧点求和即得度数公式。
最后将结论 \ref{con:xi-terminal-zm-delta-s5-holomorphic-oneforms-chevalley-weil} 中列出的不变维数代入即可得到表中数值。
证毕。
\end{proof}

\begin{remark}[关于表述的可迁移性]\label{rem:xi-terminal-zm-delta-s5-pluricanonical-transferability}
以上关于 \(V_n=H^0(X,nK_X)\) 的特征函数与不可约分解，完全由护照 \((5)\cdot(2)^6\)、切本征值谱（定理 \ref{thm:xi-terminal-zm-delta-s5-5cycle-tangent-spectrum}）以及全纯 Lefschetz 固定点公式与特征正交性推出；
因此可作为同一类 \(S_5\)-正则闭包覆盖在多重典范层面的可复用计算模板。
\end{remark}

\endinput
